\section{Conditional Probability and Independence}
\paragraph{Conditional Probability}
Probability \(E\) occurs given that \(F\) occurred is denoted \(P(E|F)\)
\[P(E|F)=\frac{P(EF)}{(F)}, P(F)>0\]
Conditional probability is a probability because it satisfies Kolmogorov's
three axioms.
\begin{itemize}
  \item \( 0\le P(E|F) \le 1 \)
  \item \( P(S|F) = 1 \)
  \item If \( E_i \) for \( i = 1,2,… \) are mutually exclusive events, then
    \[ P(\cup^\infty_{i=1}E_i|F) = \sum_{i=1}^{\infty}P(E_i|F) \]
\end{itemize}
Given that $F$ occurred, occurrence of $E$ is in $EF$. Only need to 
focus on $F$, the \emph{reduced sample space}.
\begin{itemize}
  \item \( P(EF) = P(F) P(E|F) \)
  \item \( P(E_{1}E_{2}…E_{n}) = P(E_{1}) P(E_{2}|E_{1})…P(E_n|E_{1}…E_{n-1})\)
  \item For $F_1,F_2,…,F_n$ partitioning $S$,
    \[ P(E) = \sum_{i=1}^{n} P(EF_i) = \sum_{i=1}^{n} P(F_i) P(E|F_i) \] 
  \item \textcolor{Blue}{Baye's Theorem}: $E$ occurred. Which $F_j$ most likely led to $E$:
    \begin{equation*}
      \begin{split}
        P(F_j|E) & = \frac{P(EF_j)}{P(E)} = \frac{P(F_j)P(E|F_j)}{P(E)}\\
                 & = \frac{P(F_j)P(E|F_j)}{\sum_{i=1}^{n} P(F_i)P(E|F_i)}
      \end{split}
    \end{equation*}
  \item \textcolor{Blue}{Odds} of $E=\frac{P(F)}{P(F^C)}$
\end{itemize}
\paragraph{Independent Events}
\begin{itemize}
  \item \(E\) is independent of
    $\begin{aligned}[t]
      F&\iff F \text{\ is independent of } E\\
       &\iff P(E|F) = P(E)\\
       &\iff P(EF)=P(E)P(F)\\
       &\iff E\text{\ and }F^c\text{\ are independent}
    \end{aligned}$
  \item $n$ events $E_1,E_2,\ldots,E_n$ are independent $\iff$ Any subset of
    the events are independent
\end{itemize}
